% Level 3, Week 6 -- Review and Checkpoint (student workbook)
\chapterweek{6}{Review and Checkpoint}{Connect the pieces, and use the debugger to watch your code think, one line at a time.}

\section{What you will learn}
\begin{checklist}
  \item Use the Thonny debugger to run a program one line at a time and watch the variables.
  \item Find and fix bugs in classes, sorts and searches.
  \item Show what you know in the mid-level checkpoint.
\end{checklist}

\section{The Thonny debugger}
A \textbf{debugger} is a slow-motion replay for your code. You can pause on every line and see the value of every variable.
\begin{center}
%% alt: Six Thonny buttons for debugging. Run runs at full speed. Debug starts slow motion (the bug icon). Step over runs this line and stops at the next. Step into goes inside a function call. Resume runs to the end or to a breakpoint. Stop ends the program.
\begin{tikzpicture}[btn/.style={draw=black,line width=0.8pt,rounded corners=3pt,minimum width=2.05cm,minimum height=0.85cm,font=\small\bfseries,align=center},
                    lab/.style={font=\small,align=center,text width=2.1cm}]
  \foreach \x/\name/\col/\what in {0/Run/kcgreen/{runs at full speed},
                                   2.35/Debug/kcorange/{starts slow motion (bug icon)},
                                   4.7/Step over/kcsky/{run this line, stop at the next},
                                   7.05/Step into/kcsky/{go inside a function call},
                                   9.4/Resume/kcyellow/{run to the end or a breakpoint},
                                   11.75/Stop/kcvermillion/{end the program}}{
    \node[btn,fill=\col!35] at (\x,0) {\name};
    \node[lab,below] at (\x,-0.55) {\what};
  }
\end{tikzpicture}
\end{center}
\begin{itemize}[itemsep=2pt]
  \item Turn on \menu{View > Variables}. To see an object's attributes, also turn on \menu{View > Object inspector}, then click the object in the Variables view.
  \item The highlighted line is the one that will run \emph{next}.
  \item \textbf{Breakpoint:} click a line number to put a red dot on it. Press Debug, then Resume: the program runs and pauses at the dot.
\end{itemize}

\begin{workedexample}[The debugger, click by click]
Try this with \texttt{week06/ex-bug-missing-self/starter.py}.
\begin{enumerate}[itemsep=1pt]
  \item Turn on \menu{View > Variables}. A Variables panel appears.
  \item Click the \textbf{Debug} button (the bug, next to the green Run button). The first line of the program lights up.
  \item Click \textbf{Step over}. The light moves to the next line that will run.
  \item Keep clicking Step over. Watch new names appear in the Variables panel, like \texttt{keeper}.
  \item If the program stops with a red error, read its last line: you found a bug! Fix it, then start again at step 2.
  \item When a line calls one of your methods, like \texttt{keeper.undo\_last()}, click \textbf{Step into} instead. Now you are inside \texttt{undo\_last}, line by line.
  \item When you have seen enough, click \textbf{Stop} (the red sign).
\end{enumerate}
\end{workedexample}

\section{A bug-hunting strategy}
\begin{center}
%% alt: A bug-fixing loop of five steps. 1 Read the last line of the error. 2 Predict: I think line … 3 Watch: debug and step. 4 Fix ONE thing. 5 Check: run the tests. Still failing? Go round again.
\begin{tikzpicture}[stp/.style={draw=black,line width=0.8pt,rounded corners=4pt,fill=kcsky!20,minimum height=1.1cm,text width=2.3cm,align=center,font=\small}]
  \node[stp] (a) at (0,0) {\textbf{1. Read}\\the last line of the error};
  \node[stp] (b) at (3.1,0) {\textbf{2. Predict}\\\enquote{I think line \ldots}};
  \node[stp] (c) at (6.2,0) {\textbf{3. Watch}\\debug and step};
  \node[stp] (d) at (9.3,0) {\textbf{4. Fix}\\ONE thing};
  \node[stp,fill=kcgreen!25] (e) at (12.4,0) {\textbf{5. Check}\\run the tests};
  \foreach \p/\q in {a/b,b/c,c/d,d/e} \draw[-{Stealth[length=3mm]},line width=1pt] (\p) -- (\q);
  \draw[-{Stealth[length=3mm]},line width=1pt,dashed] (e.south) to[bend left=14] node[below,font=\small]{still failing? go round again} (a.south);
\end{tikzpicture}
\end{center}

\begin{tip}[Reading an error]
Start at the \textbf{bottom}. The last line says \emph{what} went wrong (\texttt{TypeError}, \texttt{NameError}, \texttt{IndexError}\ldots). The lines right above it say \emph{where}: the line number and the line of code.
\end{tip}

\section{Weeks 1--5 in one page}
\begin{center}
\renewcommand{\arraystretch}{1.25}
\begin{tabular}{@{}p{3cm}p{10.3cm}@{}}
\toprule
\textbf{Week} & \textbf{Remember} \\ \midrule
1 Git & \texttt{git add} then \texttt{git commit -m "Say what changed"}; \texttt{log}, \texttt{diff}, \texttt{restore}. \\
2 Objects & \texttt{class}, \texttt{\_\_init\_\_}, attributes \texttt{self.x}, methods \texttt{def name(self):}. \\
3 Families & \texttt{class Child(Parent):}, \texttt{super().\_\_init\_\_(\ldots)}, overriding; \enquote{has a} = composition. \\
4 Searching & linear: one by one; binary: sorted list, check the middle, halve. \\
5 Sorting & bubble and selection sort; \texttt{sorted(key=...)}; stack (undo) and queue (turns). \\
\bottomrule
\end{tabular}
\end{center}

\begin{exercise}[Bug Hunt 1: Score Keeper]
Folder: \texttt{week06/ex-bug-missing-self}. Three bugs, all the same kind. We debug this one together. Write each error and fix:
\answerlines{3}
\textbf{Done looks like:} Run prints \texttt{Comet has 10 points}, and \texttt{check\_my\_work.py} passes.
\end{exercise}

\begin{exercise}[Bug Hunt 2: Lost Scores]
Folder: \texttt{week06/ex-bug-lost-items}. A sort that loses scores, and even empties the original list! Watch \texttt{scores}, \texttt{remaining} and \texttt{result} in the Variables view. Three bugs.\\
\textbf{Done looks like:} Run prints \texttt{Sorted:\ \ \ [10, 10, 30, 40, 50]} and \texttt{Original: [40, 10, 30, 10, 50]}.
\end{exercise}

\begin{exercise}[Bug Hunt 3: Word Finder]
Folder: \texttt{week06/ex-bug-binary-search}. Binary search with two off-by-one bugs. First trace the \emph{buggy} code by hand for \texttt{find\_word(["ant", "bee", "cat"], "ant")}:
\begin{center}
\renewcommand{\arraystretch}{1.4}
\begin{tabular}{|c|c|c|c|p{5.2cm}|}
\hline
\textbf{low} & \textbf{high} & \textbf{low < high?} & \textbf{mid} & \textbf{words[mid] and so\ldots} \\ \hline
\rule{1cm}{0pt} & \rule{1cm}{0pt} & \rule{1.8cm}{0pt} & \rule{1cm}{0pt} & \\ \hline
 & & & & \\ \hline
\end{tabular}
\end{center}
Why is \enquote{ant} never found? \rule{7cm}{0.4pt}\\
\textbf{Done looks like:} Run prints \texttt{apple -> 0}, \texttt{dragon -> 3}, \texttt{lemon -> 11} and \texttt{zebra -> -1} (\texttt{-1} means \enquote{not found}).
\end{exercise}

\begin{exercise}[Bug Hunt 4: Shape Party]
Folder: \texttt{week06/ex-bug-shape-party}. Three inheritance bugs that appear one after another: nothing drawn, then a crash, then diamonds that are only dots.\\
\textbf{Done looks like:} 4 hexagons on top and 3 diamonds (squares standing on a corner) below, 7 shapes in 7 colors.
\end{exercise}

\begin{vocab}
  \vocabword{debugger}{A tool that runs your program in slow motion so you can watch every line and variable.}
  \vocabword{step over}{Run the highlighted line in one go, then pause.}
  \vocabword{step into}{Go inside the function or method that the line calls.}
  \vocabword{breakpoint}{A marked line where the debugger pauses.}
  \vocabword{traceback}{The error report Python prints. Read it from the bottom.}
  \vocabword{trace table}{A paper table where you write each variable's value, step by step.}
\end{vocab}

\begin{reflection}
My hardest bug today was \ldots{} and what helped me find it was \ldots
\answerlines{2}
The topic I feel strongest in, and the one I want to practice:
\answerlines{2}
\end{reflection}
